% =====================================================================
%  Preambolo: pacchetti, font, ambienti e stili
% =====================================================================

% ---------- Font e lingua (richiede XeLaTeX o LuaLaTeX) ----------
\usepackage{fontspec}
\usepackage{amsmath,amssymb}
\usepackage{amsthm}
\usepackage{unicode-math}
\usepackage[italian]{babel}
\usepackage{csquotes}
\addto\captionsitalian{\renewcommand{\abstractname}{Abstract}}

% Cambria Math per gli ambienti matematici; fallback se non installato
\IfFontExistsTF{Cambria Math}
  {\setmathfont{Cambria Math}}
  {\setmathfont{Latin Modern Math}}
% Per avere anche il testo in Cambria, decommentare:
% \IfFontExistsTF{Cambria}{\setmainfont{Cambria}}{}

% ---------- Impaginazione ----------
\usepackage[a4paper,margin=2.5cm]{geometry}
\usepackage{microtype}
\usepackage{booktabs,tabularx,array,longtable}
\usepackage{enumitem}
\usepackage[font=small,labelfont=bf]{caption}
\usepackage{xcolor}
\usepackage{listings}
\usepackage{tikz}
\usetikzlibrary{arrows.meta,positioning,fit,calc,shapes.geometric,backgrounds}
\usepackage{authblk}

% ---------- Bibliografia ----------
\usepackage[backend=biber,style=numeric-comp,sorting=none,maxbibnames=6]{biblatex}
\addbibresource{bibliography/references.bib}

\usepackage[hidelinks]{hyperref}
\usepackage{xurl}

% ogni sezione inizia su una pagina nuova
\usepackage{etoolbox}
\pretocmd{\section}{\clearpage}{}{}

\setlist{itemsep=2pt,topsep=4pt}
\setlength{\parskip}{3pt}
\setlength{\emergencystretch}{2em}
\newcolumntype{L}{>{\raggedright\arraybackslash}X}
% a capo consentito dopo il trattino basso negli identificatori \texttt
\AtBeginDocument{\renewcommand{\_}{\textunderscore\allowbreak}}

% ---------- Ambiente "definition" ----------
\theoremstyle{definition}
\newtheorem{definition}{Definizione}[section]

% ---------- Listati (YANG, alberi pyang, JSON, HTTP) ----------
\definecolor{codebg}{gray}{0.96}
\definecolor{codekw}{RGB}{0,70,140}
\definecolor{codecm}{RGB}{110,110,110}
\definecolor{codestr}{RGB}{150,40,40}
\lstdefinelanguage{yang}{
  morekeywords={module,namespace,prefix,import,organization,contact,
    description,revision,reference,yang-version,identity,base,typedef,type,
    identityref,grouping,uses,container,list,key,leaf,leaf-list,config,
    augment,rpc,input,output,notification,leafref,path,default,units,
    must,when,enum,enumeration,mandatory},
  sensitive=true,
  morecomment=[l]{//},
  morecomment=[s]{/*}{*/},
  morestring=[b]"
}
\lstdefinestyle{code}{
  basicstyle=\ttfamily\footnotesize,
  keywordstyle=\color{codekw}\bfseries,
  commentstyle=\color{codecm}\itshape,
  stringstyle=\color{codestr},
  backgroundcolor=\color{codebg},
  frame=single, framerule=0pt,
  columns=fullflexible, keepspaces=true,
  xleftmargin=4pt, xrightmargin=4pt,
  breaklines=true, upquote=true,
  aboveskip=6pt, belowskip=6pt
}
\lstset{style=code}
\renewcommand{\lstlistingname}{Listato}

% ---------- Stili TikZ comuni ----------
\tikzset{
  >=Stealth,
  box/.style={draw,rounded corners,minimum height=8mm,align=center,
              font=\small},
  lbl/.style={font=\scriptsize,midway,fill=white,inner sep=1pt},
  actor/.style={draw,rounded corners,fill=gray!8,minimum width=2.4cm,
                minimum height=7mm,align=center,font=\small},
  msg/.style={->,thick},
  ret/.style={->,dashed}
}
